\documentclass[a4paper,11pt]{article}

\usepackage{revy}
\usepackage[utf8]{inputenc}
\usepackage[T1]{fontenc}
\usepackage[danish]{babel}


\revyname{MatematikRevyen}
\revyyear{2026}
% HUSK AT OPDATERE VERSIONSNUMMER
\version{0.1}
\eta{$2$ minutter}
\status{Ej færdig}

\title{Den lille gumbas}
\author{Julius Pallesen '24}

\begin{document}
\maketitle

\begin{roles}
\role{H} Henrik Holm
\end{roles}


\begin{sketch}

\says{H} Du skal ud på en rejse. En rejse ind i dit sind. En rejse helt ind i sindet. Når rejsen er forbi vil du høre mig sige nomensas, nomensas, og du vil være parat til endnu et akt\footnote{eller endnu en sketch/sang} med fornyet energi i krop og sjæl. Sørg for du sidder godt tilrette med fødderne solidt plantet i gulvet i et auditorie du holder af og omgivet af mennesker, du kan holde ud. Så er du klar til dagens drømmerejse.

\says{H} Forestil dig at du befinder dig i Vandrehallen på H.C. Ørsted-instituttet. Du står placeret nøjagtigt i midten mellem kantinen og mosaikvæg foran Auditorie 1. Du er veludhvilet, føler dig godt tilpas og din hud er ren. Du beskuer mosaikvæggen, men i periferien af dit blikfang bemærker du fire filterkaffebeholdere, to fade med glaserede sandkager og en kurv med frugt. En konference finder sted, og du mærker en pludselig sult i din vom og overvejer at tage en appelsin fra frugtkurven.

\says{H} I det du bevæger dig ned ad trappen forbi handicapelevatoren, læser du på et skilt, at konferencens emne er kontinuumshypotesen. Kontinuumshypotesen ligger langt fra dit interesseområde. Du har aldrig stået i en situation, hvor den har betydet noget som helst! Du giver nu dine tanker fri omfart i den store Vandrehal. Du kommer til at tænke på udvalgsaksiomet, og din krop og dit sind overtages af en vrede.

\says{H} Men vi skal længere ind. Vi skal helt ind i sindet. Du er nu nået hen til frugtkurven, hvor du har samlet en appelsin op. Mærk kontaktfladen mellem appelsinskrald og din hud. Du vurderer at dens volumen er cirka 270 kubikcentimeter med hensyn til øjemålet. Du skrælder nu appelsinen i seks nøje udvalgte stykker. Og pludselig bliver du \textit{suget} ind i en verden hvor volumen slet ikke giver nogen mening. Vandrehallen er væk og du står i stedet placeret på et skakbræt, og er på højde med den hvide springer ved din side. Nu er du bange.

\says{H} Men vi skal længere ind. Vi skal helt ind i sindet. Du ser nu at der ligger hundrede-, nej tusindvis af appelsiner på skakbrættets felter. Du kigger ned og opdager at du er fuldstændig splitterravende nøgen. For enden af skakbrættet for du øje på tre døre, og desperat forsøger du at mase dig gennem appelsiner. Du svømmer nærmest igennem appelsinjuice. Mærk kontaktfladen mellem smattet appelsinjuice og din nøgne hud. De sorte skakbrikker på den anden side viser sig slet ikke at være skakbrikker, men 16 kloner af din matematiklærer fra folkeskolen. Nu er du bange.

\says{H} Men vi skal længere ind. Vi skal helt ind i sindet. I det du kommer frem, hilser du venligt på de 16 folkeskolelærere, men det skulle du aldrig have gjort. De begynder at skælde dig ud over, hvor dårlig du var til at løse differentialligninger, da du gik i 6. klasse. Det er møg uretfærdigt, tænker du. Nu vil du tage flugten ud ad én af dørene. Du vælger den til højre, men bag den findes ingen flugt vej. I stedet står der en ged. En stor fed behåret ged med et langt skæg og horn i hovedet. Den begynder at løbe mod dig. Nu bliver du endnu mere bange.

\says{H} Men vi skal længere ind. Vi skal helt ind i sindet. Du glider i en appelsin og lander på ryggen. Geden står nu over dig og stirrer direkte ind i dine pupiller. Den slikker dit ansigt med sin store fede våde tunge. Mærk kontaktfladen mellem våd gedetunge og dit ansigt. Den begynder at tale til dig. "Her er hvad der sker, når man stjæler fra andre," siger geden. "Hvor våger du at tage en appelsin fra en kurv der ikke er din, dit privilegieblinde akademikerklasse københavnersvin! Frugten er kun til konferencedeltagerne!" Nu er du for alvor bange og du er \textit{pisket} til at lægge appelsinen tilbage i frugtkurven ved mosaikken foran auditoruim 1.

\says{H} Nomensas, nomensas. Her slutter dagens drømmerejse. Og du er nu parat til endnu et akt med fornyet energi i krop og sjæl.


\end{sketch}


\end{document}
